\begin{longtable}{|>{\raggedright\arraybackslash}p{3.5cm}|>{\raggedright\arraybackslash}p{3.5cm}|>{\raggedright\arraybackslash}p{4.5cm}|>{\raggedright\arraybackslash}p{6cm}|}
\caption{API Test Cases}\label{tab:api_test_cases}\\
\hline
\textbf{Test Case ID} & \textbf{Objective} & \textbf{API Endpoint \& Method} & \textbf{Request Body / Parameters \& Expected Results} \\
\hline
\endfirsthead
\multicolumn{4}{c}%
{{\bfseries \tablename\ \thetable{} -- continued from previous page}} \\
\hline
\textbf{Test Case ID} & \textbf{Objective} & \textbf{API Endpoint \& Method} & \textbf{Request Body / Parameters \& Expected Results} \\
\hline
\endhead
\hline \multicolumn{4}{|r|}{{Continued on next page}} \\ \hline
\endfoot
\hline
\endlastfoot

% --- Authentication API ---
\multicolumn{4}{|l|}{\textbf{Sub-Section: Authentication API (`/api/auth`)}} \\ \hline
TC\_API\_AUTH\_001 & Successful Librarian Login & \textbf{POST} \newline `/api/auth/login`  & \textbf{Request Body (JSON):}
\begin{lstlisting}
{
  "email": "librarian@example.com",
  "password": "valid_password"
}
\end{lstlisting}
\textbf{Expected Result:}
\begin{itemize}[noitemsep, topsep=0pt, leftmargin=*]
    \item \textbf{Status Code:} 200 OK
    \item \textbf{Response Body:} Contains user object with `isAdmin: true` and an authentication token/session cookie. 
\end{itemize} \\ \hline

TC\_API\_AUTH\_002 & Successful Member Login & \textbf{POST} \newline `/api/auth/login`  & \textbf{Request Body (JSON):}
\begin{lstlisting}
{
  "email": "member@example.com",
  "password": "valid_password"
}
\end{lstlisting}
\textbf{Expected Result:}
\begin{itemize}[noitemsep, topsep=0pt, leftmargin=*]
    \item \textbf{Status Code:} 200 OK
    \item \textbf{Response Body:} Contains user object with `isAdmin: false` and an authentication token/session cookie. 
\end{itemize} \\ \hline

TC\_API\_AUTH\_003 & Login with invalid password & \textbf{POST} \newline `/api/auth/login`  & \textbf{Request Body (JSON):}
\begin{lstlisting}
{
  "email": "librarian@example.com",
  "password": "invalid_password"
}
\end{lstlisting}
\textbf{Expected Result:}
\begin{itemize}[noitemsep, topsep=0pt, leftmargin=*]
    \item \textbf{Status Code:} 401 Unauthorized (or similar)
    \item \textbf{Response Body:} Contains an appropriate error message. 
\end{itemize} \\ \hline

TC\_API\_AUTH\_004 & Librarian registers a new member & \textbf{POST} \newline `/api/auth/register`  \newline \textit{(Requires Librarian auth)} & \textbf{Request Body (JSON):}
\begin{lstlisting}
{
  "name": "New Member",
  "email": "new.member@example.com",
  "password": "strong_password",
  "isAdmin": false
}
\end{lstlisting}
\textbf{Expected Result:}
\begin{itemize}[noitemsep, topsep=0pt, leftmargin=*]
    \item \textbf{Status Code:} 201 Created
    \item \textbf{Response Body:} Confirmation message. The user record is created in the database with a hashed password. 
\end{itemize} \\ \hline

TC\_API\_AUTH\_005 & Attempt to register with an existing email & \textbf{POST} \newline `/api/auth/register`  \newline \textit{(Requires Librarian auth)} & \textbf{Request Body (JSON):}
\begin{lstlisting}
{
  "name": "Another Member",
  "email": "member@example.com",
  "password": "another_password",
  "isAdmin": false
}
\end{lstlisting}
\textbf{Expected Result:}
\begin{itemize}[noitemsep, topsep=0pt, leftmargin=*]
    \item \textbf{Status Code:} 400 Bad Request (or similar)
    \item \textbf{Response Body:} Error message indicating the user already exists. 
\end{itemize} \\ \hline

TC\_API\_AUTH\_006 & Member attempts to access user registration endpoint & \textbf{POST} \newline `/api/auth/register`  \newline \textit{(Requires Member auth)} & \textbf{Request Body (JSON):} (any valid user data) \newline
\textbf{Expected Result:}
\begin{itemize}[noitemsep, topsep=0pt, leftmargin=*]
    \item \textbf{Status Code:} 403 Forbidden
    \item \textbf{Response Body:} Error message indicating insufficient permissions, as only librarians can register users. 
\end{itemize} \\ \hline

TC\_API\_AUTH\_007 & Successful Logout & \textbf{GET} \newline `/api/auth/logout`  \newline \textit{(Requires any user auth)} & \textbf{Request Body:} N/A \newline
\textbf{Expected Result:}
\begin{itemize}[noitemsep, topsep=0pt, leftmargin=*]
    \item \textbf{Status Code:} 200 OK (or redirect)
    \item The user's session is terminated. Subsequent requests to protected endpoints should fail. 
\end{itemize} \\ \hline

% --- Book Management API ---
\multicolumn{4}{|l|}{\textbf{Sub-Section: Book Management API (`/api/book`)}} \\ \hline
TC\_API\_BOOK\_001 & Librarian adds a new book & \textbf{POST} \newline `/api/book/add`  \newline \textit{(Requires Librarian auth)} & \textbf{Request Body (JSON):}
\begin{lstlisting}
{
  "name": "The Art of Software Testing",
  "isbn": "978-1118031534",
  "isAvailable": true
}
\end{lstlisting}
\textbf{Expected Result:}
\begin{itemize}[noitemsep, topsep=0pt, leftmargin=*]
    \item \textbf{Status Code:} 201 Created
    \item \textbf{Response Body:} The newly created book object. 
\end{itemize} \\ \hline

TC\_API\_BOOK\_002 & Member attempts to add a new book & \textbf{POST} \newline `/api/book/add`  \newline \textit{(Requires Member auth)} & \textbf{Request Body (JSON):} (any valid book data) \newline
\textbf{Expected Result:}
\begin{itemize}[noitemsep, topsep=0pt, leftmargin=*]
    \item \textbf{Status Code:} 403 Forbidden
    \item \textbf{Response Body:} Error message indicating insufficient permissions. Librarians have exclusive add rights. 
\end{itemize} \\ \hline

TC\_API\_BOOK\_003 & Any user gets a list of all books & \textbf{GET} \newline `/api/book/getAll`  \newline \textit{(Public or Member auth)} & \textbf{Request Body:} N/A \newline
\textbf{Expected Result:}
\begin{itemize}[noitemsep, topsep=0pt, leftmargin=*]
    \item \textbf{Status Code:} 200 OK
    \item \textbf{Response Body:} An array of book objects. 
\end{itemize} \\ \hline

TC\_API\_BOOK\_004 & Librarian deletes a book & \textbf{DELETE} \newline `/api/book/delete/{id}`  \newline \textit{(Requires Librarian auth)} & \textbf{Parameters:} URL parameter `id` of an existing book. \newline
\textbf{Expected Result:}
\begin{itemize}[noitemsep, topsep=0pt, leftmargin=*]
    \item \textbf{Status Code:} 200 OK (or 204 No Content)
    \item \textbf{Response Body:} Success message. The book is removed from the database.
\end{itemize} \\ \hline

TC\_API\_BOOK\_005 & Unauthenticated user attempts to delete a book & \textbf{DELETE} \newline `/api/book/delete/{id}`  \newline \textit{(No auth)} & \textbf{Parameters:} URL parameter `id` of an existing book. \newline
\textbf{Expected Result:}
\begin{itemize}[noitemsep, topsep=0pt, leftmargin=*]
    \item \textbf{Status Code:} 401 Unauthorized
    \item \textbf{Response Body:} Error message indicating authentication is required. 
\end{itemize} \\ \hline

% --- Borrowal Management API ---
\multicolumn{4}{|l|}{\textbf{Sub-Section: Borrowal Management API (`/api/borrowal`)}} \\ \hline
TC\_API\_BORROW\_001 & Member borrows an available book for themselves & \textbf{POST} \newline `/api/borrowal/add`  \newline \textit{(Requires Member auth)} & \textbf{Request Body (JSON):}
\begin{lstlisting}
{
  "bookId": "valid_book_id",
  "memberId": "current_user_id"
}
\end{lstlisting}
\textbf{Expected Result:}
\begin{itemize}[noitemsep, topsep=0pt, leftmargin=*]
    \item \textbf{Status Code:} 201 Created
    \item \textbf{Response Body:} The new borrowal record. The book's availability may be updated. 
\end{itemize} \\ \hline

TC\_API\_BORROW\_002 & Member attempts to borrow a book for another member & \textbf{POST} \newline `/api/borrowal/add`  \newline \textit{(Requires Member auth)} & \textbf{Request Body (JSON):}
\begin{lstlisting}
{
  "bookId": "valid_book_id",
  "memberId": "another_user_id"
}
\end{lstlisting}
\textbf{Expected Result:}
\begin{itemize}[noitemsep, topsep=0pt, leftmargin=*]
    \item \textbf{Status Code:} 403 Forbidden
    \item \textbf{Response Body:} Error message. The API must enforce that the `memberId` in the request matches the authenticated user's ID. 
\end{itemize} \\ \hline

TC\_API\_BORROW\_003 & Librarian updates a borrowal status (e.g., to 'Returned') & \textbf{PUT} \newline `/api/borrowal/update/{id}`  \newline \textit{(Requires Librarian auth)} & \textbf{Parameters:} URL parameter `id` of an existing borrowal. \newline
\textbf{Request Body (JSON):}
\begin{lstlisting}
{
  "status": "Returned"
}
\end{lstlisting}
\textbf{Expected Result:}
\begin{itemize}[noitemsep, topsep=0pt, leftmargin=*]
    \item \textbf{Status Code:} 200 OK
    \item \textbf{Response Body:} The updated borrowal object.
\end{itemize} \\ \hline

TC\_API\_BORROW\_004 & Member attempts to view all borrowals in the system & \textbf{GET} \newline `/api/borrowal/getAll`  \newline \textit{(Requires Member auth)} & \textbf{Request Body:} N/A \newline
\textbf{Expected Result:}
\begin{itemize}[noitemsep, topsep=0pt, leftmargin=*]
    \item \textbf{Status Code:} 200 OK
    \item \textbf{Response Body:} An array containing \textbf{only} the borrowal records for the authenticated member. The API should not leak other users' data. 
\end{itemize} \\ \hline

% --- Security Testing ---
\multicolumn{4}{|l|}{\textbf{Sub-Section: Security \& Input Validation}} \\ \hline
TC\_API\_SEC\_001 & Verify password hash is not returned on user lookup & \textbf{GET} \newline `/api/user/get/{id}`  \newline \textit{(Requires Librarian auth)} & \textbf{Parameters:} URL parameter `id` for any user. \newline
\textbf{Expected Result:}
\begin{itemize}[noitemsep, topsep=0pt, leftmargin=*]
    \item \textbf{Status Code:} 200 OK
    \item \textbf{Response Body:} The user object \textbf{must not} contain the `hash` or `salt` fields. 
\end{itemize} \\ \hline

TC\_API\_SEC\_002 & Attempt basic SQL Injection (or NoSQL equivalent) & \textbf{POST} \newline `/api/auth/login`  & \textbf{Request Body (JSON):}
\begin{lstlisting}
{
  "email": "' OR '1'='1",
  "password": "' OR '1'='1"
}
\end{lstlisting}
\textbf{Expected Result:}
\begin{itemize}[noitemsep, topsep=0pt, leftmargin=*]
    \item \textbf{Status Code:} 401 Unauthorized (or 400 Bad Request)
    \item \textbf{Response Body:} Normal login failure message. The system must not be compromised. This validates server-side input validation. 
\end{itemize} \\ \hline

\end{longtable}